\documentclass[12pt]{article}
\usepackage{amsmath,amssymb,amsfonts,mathrsfs}
\usepackage[margin=1in]{geometry}

\begin{document}

\begin{center}
{\Large\bfseries Chapter 10, Problem 17}
\end{center}

\noindent\textbf{Problem.} Using the three-dimensional representation of the rotation group, given, for example, by Eq.~(1.25), find the matrices for a three-dimensional representation of the symmetry group, $\mathscr{O}$, of the cube. We saw in Section 10.8 that there were two three-dimensional representations of $\mathscr{O}$. Give a simple prescription for finding the matrices belonging to the second three-dimensional representation (i.e., the one not given by the rotation group).

\bigskip
\noindent\textbf{Solution.}

The rotation group of the cube, $\mathscr{O}$, has 24 elements. Using the three-dimensional representation given by $3\times 3$ orthogonal rotation matrices acting on $(x, y, z)$:

\medskip
\noindent\textbf{The 24 rotation matrices:}

We organize by class (using the character table from Section 10.8):

\textbf{Class 1}: Identity (1 element)
\[
D(e) = \begin{pmatrix} 1 & 0 & 0 \\ 0 & 1 & 0 \\ 0 & 0 & 1 \end{pmatrix}.
\]

\textbf{Class 2}: Six $90^\circ$ rotations about coordinate axes (6 elements).

About $z$-axis by $\pm 90^\circ$:
\[
R_z(90^\circ) = \begin{pmatrix} 0 & -1 & 0 \\ 1 & 0 & 0 \\ 0 & 0 & 1 \end{pmatrix}, \quad
R_z(-90^\circ) = \begin{pmatrix} 0 & 1 & 0 \\ -1 & 0 & 0 \\ 0 & 0 & 1 \end{pmatrix}.
\]
Similarly for $x$ and $y$ axes (4 more matrices by cyclic permutation of coordinates).

\textbf{Class 3}: Three $180^\circ$ rotations about coordinate axes (3 elements).
\[
R_z(180^\circ) = \begin{pmatrix} -1 & 0 & 0 \\ 0 & -1 & 0 \\ 0 & 0 & 1 \end{pmatrix},
\]
and similarly $R_x(180^\circ)$, $R_y(180^\circ)$.

\textbf{Class 4}: Eight $120^\circ$ rotations about body diagonals (8 elements).

For the diagonal $(1,1,1)$, rotation by $120^\circ$ cyclically permutes $x \to y \to z \to x$:
\[
R_{111}(120^\circ) = \begin{pmatrix} 0 & 0 & 1 \\ 1 & 0 & 0 \\ 0 & 1 & 0 \end{pmatrix}, \quad
R_{111}(240^\circ) = \begin{pmatrix} 0 & 1 & 0 \\ 0 & 0 & 1 \\ 1 & 0 & 0 \end{pmatrix}.
\]
Similarly for the other three diagonals $(1,1,-1)$, $(1,-1,1)$, $(-1,1,1)$, giving 6 more matrices.

\textbf{Class 5}: Six $180^\circ$ rotations about face-diagonal axes (6 elements).

About the axis bisecting the $x$ and $y$ axes:
\[
R_{xy}(180^\circ) = \begin{pmatrix} 0 & 1 & 0 \\ 1 & 0 & 0 \\ 0 & 0 & -1 \end{pmatrix},
\]
and similarly for the other five face-diagonal axes.

\medskip
\noindent\textbf{Characters of this representation:}

\[
\begin{array}{c|ccccc}
\text{Class} & 1 & 2 & 3 & 4 & 5 \\
g_i & 1 & 6 & 3 & 8 & 6 \\
\hline
\chi & 3 & 1 & -1 & 0 & -1
\end{array}
\]

This matches one of the two 3D irreducible representations from the character table (10.91).

\medskip
\noindent\textbf{Finding the second 3D representation:}

The second three-dimensional representation can be obtained by a simple prescription: \textbf{multiply each rotation matrix by the determinant of the corresponding permutation character}, or more directly, multiply by the characters of the one-dimensional ``alternating'' representation.

Specifically, from the character table of $\mathscr{O}$, the second 3D irreducible representation has characters:
\[
\chi' = (3, -1, -1, 0, 1).
\]

The relationship between the two is: $\chi'_i = \chi^{(\Gamma_2)}_i \cdot \chi_i$, where $\Gamma_2$ is the nontrivial 1D representation with characters $(1, -1, 1, 1, -1)$.

Therefore, the matrices $D'(g)$ of the second 3D representation are obtained from the first by:
\[
\boxed{D'(g) = \chi^{(\Gamma_2)}(g) \cdot D(g),}
\]
i.e., multiply each matrix by $+1$ or $-1$ according to the alternating representation. This means the $90^\circ$ rotations and $180^\circ$ face-diagonal rotations get a sign flip, while the identity, $180^\circ$ coordinate-axis rotations, and $120^\circ$ diagonal rotations are unchanged.

\end{document}
